
%%%%%%%%%%%%%%%%% 证明方法 %%%%%%%%%%%%%%

\chapter{证明方法}

学会了简单的逻辑后，接下来便是学习如何证明。在本章中我们将介绍一些证明的方法。这
里我们只谈有关证明的一些基本原则，而不谈证明的技巧，所以给的例子将会挑选浅显易懂
的证明。

\section{若-则的证明}

在数学中最常看到的就是这种 $P \Rightarrow Q$ 的命题。要证明这种命题，我们大致上有直接法（direct method）\index{直接法（direct method）}、逆否命题法（contrapositive method）和反证法（contradiction method）三种方法。

\subsection{直接法}

所谓直接法，就是直接利用$P$成立的假设得到 $Q$ 成立。（再次强调，我们不必管 $P$ 不成立时，$Q$ 会如何）。当我们要证明 $P \Rightarrow Q$
时，若觉得 $P$ 的条件已经足够，便可以考虑使用直接证明。例如以下的例子.

\begin{example}
设 $p,a,b$ 为整数。若 $p \mid a$（$p$能整除$a$）且 $p \mid b$，则 $p \mid (a+b)$。
\end{example}

\begin{proof}
由假设 $p \mid a$ 且 $p \mid b$，知存在整数 $m,n$ 使得 $a = pm$, $b = pn$。故得
$$a + b = pm + pn = p(m + n).$$
因 $m + n$ 为整数，得证 $p \mid a + b$。
\end{proof}

有时用直接证明的方法并不能一次到位，需要借助其他的结果帮忙才能完成。也就是说或许我们不能直接证出 $P\Rightarrow Q$，但若证得 $P\Rightarrow R$ 又证得 $R\Rightarrow Q$，此时便证得 $P\Rightarrow Q$ 了。这是因为若证得 $P\Rightarrow R$，表示 $P$ 对的话 $R$ 一定对，再由 $R\Rightarrow Q$ 知 $R$ 对的话 $Q$ 一定对，故连结而知 $P$ 对则 $Q$ 一定对，得证 $P\Rightarrow Q$。这种利用传递性的证明通常有一部分是一些常用的性质或是一些（辅助）定理。例如以下的例子。

\begin{example}
设 $a$ 为正实数且 $a\neq 1$。已知若 $a^{z} = 1$，则 $z = 0$。证明若 $x,y$ 为实数满足 $a^{x} = a^{y}$，则 $x = y$。
\end{example}

\begin{proof}
由于 $a \neq 0$，对于任何实数 $y$，我们知道 $a^{y} \neq 0$，故由 $a^{x} = a^{y}$，等号两边除以 $a^{y}$ 得 $a^{x - y} = 1$。又因 $a \neq 1$，我们知道若 $a^{z} = 1$，则 $z = 0$。故由 $a^{x - y} = 1$ 可得 $x - y = 0$，得证 $x = y$。
\end{proof}

这个证明的例子中，我们其实是先证得 $(a^x = a^y) \Rightarrow (a^{x - y} = 1)$，再由 $(a^{x - y} = 1) \Rightarrow (x = y)$ 得证 $(a^x = a^y) \Rightarrow (x = y)$。其中我们用了一个大家都知道的事实，即当 $a$ 为正实数且 $a \neq 1$，若 $a^{z} = 1$，则 $z = 0$。这个事实是需要证明的，不过不容易用直接法证明，等一下我们会利用反证法来证明。

有时在直接法中我们可以分成好几种情况，看看哪些情况符合 $P$ 的条件，然后证得 $Q$。这样的证明方法有时称为\textbf{分类讨论（proof in cases）}。\index{分类讨论（proof in cases）}例如以下的例子。

\begin{example}
假设 $x$ 为实数。证明若 $x^2 - 3x + 2 < 0$，则 $1 < x < 2$。
\end{example}

\begin{proof}
由 $x^2 - 3x + 2 = (x - 1)(x - 2) < 0$，我们知道可分成 2 种情况，即
\[(1) (x - 1) < 0 \text{ 且 } (x - 2) > 0;\]
\[(2) (x - 1) > 0 \text{ 且 } (x - 2) < 0.\]
(1) 的情况表示 $x < 1$ 且 $x > 2$。由于没有实数 $x$ 会同时满足 $x < 1$ 以及 $x > 2$，我们知道 (1) 不可能成立，故推得 (2)，即 $x > 1$ 且 $x < 2$。证得 $1 < x < 2$。
\end{proof}

注意，在这个证明中，有些同学或许会疑惑为什么是排除 (1)，而不直接验证 (2) 可得 $x^2 - 3x + 2 < 0$ 呢？这是错误的。在我们的证明中明明白白表示：若 $x$ 满足 $x^2 - 3x + 2 < 0$，那么 $x$ 一定会满足 (1) 或是 (2)。排除 (1) 表示只有 (2) 会对，所以确定若 $x$ 满足 $x^2 - 3x + 2 < 0$，那么 $x$ 一定满足 (2)。若仅说 $x$ 满足 (2) 可得 $x^2 - 3x + 2 < 0$，而没有排除 (1)，那么这是证明若 $1 < x < 2$，则 $x^2 - 3x + 2 < 0$，而不是证若 $x^2 - 3x + 2 < 0$，则 $1 < x < 2$。千万别搞错。

\begin{question}
假设 $x$ 为实数。
\begin{enumerate}
    \item “若 $x^2 - 3x + 2 < 0$，则 $0 < x < 3$。” 和 “若 $x^2 - 3x + 2 < 0$，则 $1.3 < x < 1.7$。” 这两个命题哪一个是对的？
    \item “若 $0 < x < 3$，则 $x^2 - 3x + 2 < 0$。” 和 “若 $1.3 < x < 1.7$，则 $x^2 - 3x + 2 < 0$。” 这两个命题哪一个是对的？
\end{enumerate}
\end{question}

有时待论证的问题有多种结论，例如要论证 $P \Rightarrow Q \land R$ 或是 $P \Rightarrow Q \lor R$。当要论证 $P \Rightarrow Q \land R$，我们当然就是证明 $P$ 对则 $Q$ 和 $R$ 都会对。但要证明 $P \Rightarrow Q \lor R$，就会有点麻烦。因为 $Q \lor R$ 会对表示有可能 $Q$ 对 $R$ 错； $Q$ 错 $R$ 对，甚至 $Q, R$ 都对。也因此，要推得 $Q \lor R$ 会对，感觉有点麻烦。接下来这个技巧就很好用。我们可以直接仅考虑已知 $P$ 对且 $Q$ 错的情况，因为如果 $Q$ 对那自然 $Q \lor R$ 是对的。因此，此时我们就只要证明由 $P$ 对 $Q$ 错可推得 $R$ 对即可。从逻辑来看，因为 $(P \Rightarrow (Q \lor R)) \sim (\neg P \lor (Q \lor R))$，而 $((P \land \neg Q) \Rightarrow R) \sim (\neg (P \land \neg Q) \lor R) \sim ((\neg P \lor Q) \lor R)$，所以 $(P \Rightarrow (Q \lor R)) \sim ((P \land \neg Q) \Rightarrow R)$。我们看以下的例子。

\begin{example}
设 $x, y$ 为实数。证明若 $xy = 0$，则 $x = 0$ 或 $y = 0$。
\end{example}

\begin{proof}
我们可假设 $xy = 0$ 且 $x \ne 0$，此时因$x \ne 0$，可将等式$xy=0$的两边除以$x$（或说乘以$1/x$），得$y = 0$。
\end{proof}

当然，在证明$P\Rightarrow (Q\lor R)$时，若觉得$\neg R$比较好用，也可改成证$(P\lor \neg R)\Rightarrow Q$。

\subsection{逆否命题法}

我们称 $(\neg Q) \Rightarrow (\neg P)$ 为 $P \Rightarrow Q$ 这个命题的\textbf{逆否命题（contrapositive statement）}。\index{逆否命题（contrapositive statement）}回顾一下，我们知道 $P \Rightarrow Q$ 和 $(\neg Q) \Rightarrow (\neg P)$ 是逻辑等价的（参见式子 (1.11)）。也就是说 $P \Rightarrow Q$ 和 $(\neg Q) \Rightarrow (\neg P)$ 的对错是一致的。因此，若我们能证明 $(\neg Q) \Rightarrow (\neg P)$ 便证得 $P \Rightarrow Q$ 了。

当要证明 $P \Rightarrow Q$ 时，若发现 $P$ 的条件似乎不容易帮助我们证明，而 $\neg Q$ 较容易处理时，便可以考虑使用逆否命题法，也就是说证明 $(\neg Q) \Rightarrow (\neg P)$。最常发生的情况就是有不等式的情形。因为不等式不如等式使用方便，不等式的很多使用规则其实都是用等式推导的，所以如果一个命题牵涉到不等式，而它的逆否命题是等式，那么自然用逆否命题法会比较容易证明。我们有以下例子。

\begin{example}
设 $x, y$ 为实数。证明若 $x \neq y$，则 $x^3 \neq y^3$。
\end{example}

当然，若你了解 $f(x) = x^3$ 的图形或利用微积分，可以知道这一定是对的。不过我们想要用比较基础的方法处理。若用逆否命题法来证明，就是先假设 $\neg (x^3 \neq y^3)$（即 $x^3 = y^3$），要证得 $\neg (x \neq y)$（即 $x = y$）。

\begin{proof}
利用逆否命题法，首先假设 $x^3 = y^3$，即 $0 = x^3 - y^3 = (x - y)(x^2 + xy + y^2)$。由 例 2.1.4 可得 $x - y = 0$ 或 $x^2 + xy + y^2 = 0$。因此可用分类讨论处理。

(1) $x - y = 0$：此时即 $x = y$。

(2) $x^2 + xy + y^2 = 0$：此时由
$$x^2 + xy + y^2 = (x + \frac{1}{2}y)^2 + \frac{3}{4}y^2.$$
可得 $x + \frac{1}{2} y = 0$ 且 $y = 0$。因此得 $x = y = 0$。

由于所有情况皆得 $x = y$，故得证若 $x^3 = y^3$ 则 $x = y$，也因此得证若 $x \neq y$，则 $x^3 \neq y^3$。
\end{proof}

当我们利用逆否命题法把要证明的 $P \Rightarrow Q$ 改为证明 $\neg Q \Rightarrow \neg P$ 后，就可以用前面所提的直接法的方法证明 $\neg Q \Rightarrow \neg P$。在上例中，我们就是利用传递性以及分类讨论处理。

若一个命题是由包含 $x$ 的式子的性质，来推得 $x$ 本身的性质，那么由于通常是反过来将 $x$ 代入式子比较容易，此时也是用逆否命题法的好时机。例如以下的简单例子。

\begin{example}
设 $x$ 为整数。证明若 $x^2$ 是偶数，则 $x$ 是偶数。
\end{example}

\begin{proof}
用逆否命题法，即证明若 $x$ 为奇数，则 $x^2$ 为奇数。然而 $x$ 为奇数，表示 $x$ 可以写成 $x = 2n + 1$，其中 $n$ 为整数。故得证 $x^2 = (2n + 1)^2 = 4n^2 + 4n + 1 = 2(2n^2 + 2n) + 1$ 为奇数。
\end{proof}

\begin{question}
假设 $x, y$ 为整数。试用逆否命题法证明若 $x + y$ 是偶数，则 $x$ 和 $y$ 同为偶数或同为奇数。
\end{question}

有时用分类讨论处理的情况，也可用逆否命题法来证明。例如 例 2.1.3 就可以用逆否命题法证明。也就是说，先假设 $\neg (1 < x < 2)$（即 $x \ge 2$ 或 $x \le 1$）求得 $\neg (x^2 - 3x + 2 < 0)$（即 $x^2 - 3x + 2 \ge 0$）。我们看另外的例子。

\begin{example}
设 $x, y, a$ 为正实数。证明若 $xy = a$，则 $x \le \sqrt{a}$ 或 $y \le \sqrt{a}$。
\end{example}

用分类讨论证明，需考虑 $x, y$ 所有可能与 $\sqrt{a}$ 的大小关系，要分成好几种情况。而用逆否命题法，我们只要考虑一种情况。

\begin{proof}
用逆否命题法，先假设 $\neg ((x \le \sqrt{a}) \lor (y \le \sqrt{a}))$（即 $x > \sqrt{a}$ 且 $y > \sqrt{a}$）要证得 $\neg (xy = a)$（即 $xy \neq a$）。然而 $x, y, a$ 为正，故由 $x > \sqrt{a}$ 且 $y > \sqrt{a}$ 可得 $xy > (\sqrt{a})^{2} = a$，得证 $xy \neq a$。
\end{proof}

\begin{question}
利用类似 例 2.1.4 中证明 $P \Rightarrow (Q \lor R)$ 的方法，证明 例 2.1.7。
\end{question}

\begin{question}
设 $x, y, a$ 为正实数。请问若 $xy = a$，则 $x \ge \sqrt{a}$ 或 $y \ge \sqrt{a}$ 是否为对？若为对，这和 例 2.1.7 是否有矛盾？
\end{question}

\subsection{反证法}

回顾式子 (1.10) 告诉我们 $P \Rightarrow Q$ 和 $Q \lor \neg P$ 是逻辑等价的。也就是说，若我们能证明 $Q \lor \neg P$ 就等同于证明了 $P \Rightarrow Q$。然而 $Q \lor \neg P$ 是 “或” 的情况，处理起来有点像前面提过的分类讨论，没有太多的优势。不过若考虑 $\neg (P \Rightarrow Q)$，此时和 $(\neg Q) \land P$ 为逻辑等价（参见式子 (1.12)）。因此若能证明 $(\neg Q) \land P$ 一定是错的，便证得 $P \Rightarrow Q$ 为对。这就是所谓的反证法（contradiction method）\index{反证法（contradiction method）}。

反证法的处理方法就是，先假设 $(\neg Q)$ 和 $P$ 皆为对，再从中推导出与我们知道一定对的命题相矛盾。如此一来便表示 $(\neg Q)$ 和 $P$ 是错的，而得证 $P \Rightarrow Q$。这个方法的最大优点就是它一次便同时假设 $P$ 和 $\neg Q$ 为对，给我们较多的信息去推导，而不像直接法仅假设 $P$ 为对，或逆否命题法仅假设 $\neg Q$ 为对。它的缺点就是，不像直接法或逆否命题法明确地知道要推导什么（即由 $P$ 推导出 $Q$ 或由 $\neg Q$ 推导出 $\neg P$），而是要推导出一个 “未知” 的矛盾。

当要证明 $P \Rightarrow Q$ 时，若发现单独 $P$ 的条件或是单独 $\neg Q$ 的条件似乎不容易帮助我们证明，这时便可以考虑 $P$ 和 $\neg Q$ 的条件可同时使用的反证法。我们看以下的例子。

\begin{example}
设 $r$ 为实数，证明若 $r^{2} = 2$，则 $r$ 是无理数。
\end{example}

我们几乎没有直接的方法证明一个数是无理数，所以不可能用直接法。而若用逆否命题法虽可先假设 $r$ 为有理数，但要推得 $r^{2} \neq 2$ 这个不等式。前面已提过不等式的推导并不容易，所以我们用反证法来证明。

\begin{proof}
用反证法，即假设 $r$ 为有理数且满足 $r^{2} = 2$，希望能得到矛盾。依假设 $r$ 为有理数，表示 $r$ 可以写成 $r = m / n$，其中 $m, n$ 为整数。现若 $m, n$ 皆为偶数，我们可以约掉 2，如此一直下去，我们可假设 $m, n$ 为一奇一偶或皆为奇数。然而 $r^{2} = 2$，即 $m^{2} = 2n^{2}$ 为偶数，故由 例 2.1.6 知 $m$ 必为偶数。也就是说 $m$ 可写成 $m = 2m^{\prime}$，其中 $m^{\prime}$ 为整数。此时得 $4m^{\prime 2} = 2n^{2}$，即 $n^{2} = 2m^{\prime 2}$ 为偶数。故再由 例 2.1.6 知 $n$ 亦为偶数。此与当初假设 $m, n$ 为一奇一偶或皆为奇数相矛盾。故得证若 $r^{2} = 2$，则 $r$ 是无理数。
\end{proof}

从 例 2.1.8 中我们应可体会，当初若没有想到将 $r = m / n$ 中分子分母的 2 约干净，就无法推出矛盾了。所以用反证法证明的困难处就是要 “制造矛盾”。

回顾在 例 2.1.2 的证明中，我们用了一个事实，即当 $a \neq 1$ 且为正实数时，若 $z$ 为实数满足 $a^{z} = 1$，则 $z = 0$。我们可以用反证法来证明这个命题。

\begin{example}
设 $a \neq 1$ 且为正实数。证明若 $z$ 为实数满足 $a^{z} = 1$，则 $z = 0$。
\end{example}

若要用反证法，我们必须假设 $z \neq 0$ 且 $a^{z} = 1$ 而得到矛盾。要如何制造矛盾呢？我们知道这个命题中 $a \neq 1$ 是重要的前提（否则它会错），所以矛盾的关键是 $a \neq 1$。

\begin{proof}
我们利用反证法，先假设 $z \neq 0$ 且 $a^{z} = 1$。此时由于 $z \neq 0$，我们知道 $1 / z$ 是存在的，故利用 $(a^{z})^{1 / z} = a$ 以及 $a^{z} = 1$，得
$$a = (a^{z})^{1 / z} = 1^{1 / z} = 1.$$
此与已知 $a \neq 1$ 相矛盾，得证若 $a^{z} = 1$，则 $z = 0$。
\end{proof}

\begin{remark}
在证明 $P \Rightarrow Q$ 时，其实用逆否命题法或反证法它们第一步是一样的，也就是先假设 $\neg Q$。接着我们就可以判断条件是否够我们推导出 $\neg P$。若条件够，那就用逆否命题法。若条件不够，则可以加上 $P$ 的条件，看看是否可推得矛盾，这就是反证法。这两种方法，各有很多种中文名称，甚至两种分别都有人称之为 “反证法”。为了方便起见，以后我不太想去区分这两种方法，我都一律称之为反证法。也就说只要是要用 $\neg Q$ 这样的条件来证明，我都一律称之为 “反证法”。
\end{remark}

\subsection{当且仅当的证明}

$P \Leftrightarrow Q$ 的证明基本上就是要证明 $P \Rightarrow Q$ 和 $Q \Rightarrow P$。大家或许看过有些证明由于每一步逆推回去也是对的，所以在推导完 $P \Rightarrow Q$ 后便说反向亦然，而得证 $P \Leftrightarrow Q$。在这里特别提醒大家，除非你确认每一步逆推回去也是对的，千万不要随便认定反向亦然，就说得到 $P \Leftrightarrow Q$。尤其在以后许多进阶一点的定理，很可能两边的推导方式是用到完全不同的概念或原理。所以证明 $P \Leftrightarrow Q$ 还是要分别证明 $P \Rightarrow Q$ 和 $Q \Rightarrow P$ 为宜。

例如设 $a$ 为整数要证明 “$a^2$ 为偶数 $\Leftrightarrow a$ 为偶数”，大家可能会先证 $\Leftarrow$ 这个方向。此时可令 $a = 2n$，其中 $n$ 为整数，使得 $a^2 = (2n)^2 = 4n^2$ 为偶数。但这个推导方式，逆推就有问题了。因为知道 $a^2$ 为偶数，为何 $a^2$ 一定可以写成 $4n^2$ 这个样子呢？所以 $\Rightarrow$ 这个方向还是要推导的（我们在 例 2.1.6 已证明过了）。

由于 $(P\Rightarrow Q)\sim ((\neg Q)\Rightarrow (\neg P))$ 且 $(Q\Rightarrow P)\sim ((\neg P)\Rightarrow (\neg Q))$，我们知 $P\Leftrightarrow Q$ 和 $(\neg P)\Leftrightarrow (\neg Q)$ 为逻辑等价。有的同学会觉得若证明 $P\Leftrightarrow Q$ 较麻烦，可考虑证明 $(\neg P)\Leftrightarrow (\neg Q)$。其实这是没必要的，不管证明哪一个都需要经过一样的过程。例如假设 $a,b$ 为整数，证明 “$ab$为偶数 $\Leftrightarrow a$ 为偶数或 $b$ 为偶数” 和证明 “$ab$为奇数 $\Leftrightarrow a$ 和 $b$都为奇数” 其实是一样的，即使后面那一个看起来比较简洁，但是不管证明哪一个，证明的过程都是一样的。反倒是有的同学可能会证明了 $(\neg Q)\Rightarrow (\neg P)$ 后又去证 $P\Rightarrow Q$，重复证明了同一件事而不自知，误以为证得了两个方向，千万要注意。所以基本上在证明当且仅当的命题 时，最好表明目前在证明哪一个方向。这样自己较不会弄错，看证明的人也较清楚整个证明过程。两全其美，何乐而不为呢？

数学上也经常会有类似这样的命题：
\begin{center}
    下列等价： (1) $P$ ; (2) $Q$ ; (3) $R$。
\end{center}
（有时不只会有 $P,Q,R$ 三项，可能会有更多项）。这个意思就是
$$P\Leftrightarrow Q,\; Q\Leftrightarrow R\; \text{且}\; R\Leftrightarrow P.$$
虽然如此，我们不必把六个 $\Rightarrow$ 都证。事实上仅要证明
$$P\Rightarrow Q,\; Q\Rightarrow R\; \text{且}\; R\Rightarrow P$$
即可。这是因为 $Q\Rightarrow P$ 的部分，可由 $Q\Rightarrow R$ 以及 $R\Rightarrow P$ 得到，而 $R\Rightarrow Q$ 的部分，可由 $R\Rightarrow P$ 以及 $P\Rightarrow Q$ 得到。最后 $P\Rightarrow R$ 的部分，可由 $P\Rightarrow Q$ 以及 $Q\Rightarrow R$ 得到。当然，有时不一定这个顺序好证，也可考虑倒过来证明
$$R\Rightarrow Q,\; Q\Rightarrow P\; \text{且}\; P\Rightarrow R.$$
甚至有时候会发生不管哪一边都不好证，例如 $P\Leftrightarrow R$ 两边都不好证，这时证明
$$P\Leftrightarrow Q\; \text{且}\; Q\Leftrightarrow R$$
也可。因为 $P\Rightarrow R$ 的部分，可由 $P\Rightarrow Q$ 以及 $Q\Rightarrow R$ 得到，而 $R\Rightarrow P$ 的部分，可由 $R\Rightarrow Q$ 以及 $Q\Rightarrow P$ 得到。总之，在证明这一类的问题时，要注意推导的方向，确保任两个命题 都可通行无阻。时时标明目前是从哪一个命题 推导到哪一个命题，是必要的。

\section{存在性和唯一性的证明}

有关存在性（existence）\index{存在性（existence）}和唯一性（uniqueness）\index{唯一性（uniqueness）}的探讨也经常在数学定理中出现。要注意，存在性和唯一性是两个互相独立的性质。也就是说存在未必会唯一。而这里的唯一指的是若存在则会唯一，所以证得唯一也未必会存在。所以这里，我们分开讨论存在性和唯一性的证明。

\subsection{存在性}

有关存在性的证明方法大致上有两类。一类是所谓的构造性方法（constructive method），指的是确实告知存在的是什么。另一类是非构造性方法（nonconstructive method），指的是利用已知的理论或逻辑的推导得知一定存在，但未必知道有哪些。

例如证明存在实数 $x$ 满足 $6x^{2} - x - 1 = 0$。我们可以将 $6x^{2} - x - 1$ 分解得 $(2x - 1)(3x + 1)$ 故明确找出 $x = 1 / 2$（或 $x = - 1 / 3$）这个实数会满足 $6x^{2} - x - 1 = 0$。这就是一个构造性方法。我们也可考虑多项式函数 $f(x) = 6x^{2} - x - 1$，发现 $f(0) = - 1 < 0$ 且 $f(1) = 4 > 0$。故由多项式函数为连续函数以及连续函数的中间值定理，证得 $f(x) = 0$ 在 $0 < x < 1$ 之间必有一根，而证得了存在性。这个方法虽证出存在性，但因无法明确指出哪一个 $x$ 会满足 $6x^{2} - x - 1 = 0$，所以是非构造性方法。我们再看以下的例子。

\begin{example}
证明存在无理数 $a, b$使得 $a^{b}$为有理数。
\end{example}

\begin{proof}
\textbf{构造性方法}：考虑 $a = \sqrt{2}$ 且 $b = \log_{2}9$。我们已知 $a$ 为无理数。同样利用反证法可证明 $b$ 亦为无理数。事实上若存在 $m, n$ 为整数满足 $\log_{2}3 = n / m$，就表示 $2^{n} = 3^{m}$，我们知道 $3$ 的任何整数次方都不会是偶数，得到矛盾。故知 $b = \log_{2}9$ 为无理数。此时
$$a^{b} = 2^{\frac{1}{2}\log_{2}9} = 2^{\log_{2}3} = 3,$$
得证存在无理数 $a, b$ 使得 $a^{b}$ 为有理数。

\textbf{非构造性方法}：考虑 $a^{\prime} = \sqrt{2}$ 且 $b^{\prime} = \sqrt{2}$。我们知道 $a^{\prime}, b^{\prime}$ 为无理数。令 $c = (a^{\prime})^{b^{\prime}}$。若 $c$ 为有理数，则 $a = \sqrt{2}, b = \sqrt{2}$ 为所求。而若 $c$ 为无理数，此时令 $a = c, b = \sqrt{2}$，则
$$a^{b} = (\sqrt{2}^{\sqrt{2}})^{\sqrt{2}} = \sqrt{2}^{2} = 2.$$
得证存在无理数 $a, b$ 使得 $a^{b}$ 为有理数。
\end{proof}

这里这个非构造性方法由于没有证明 $\sqrt{2}^{\sqrt{2}}$ 是否为有理数。因此无法确定 $a = b = \sqrt{2}$ 和 $a = \sqrt{2}^{\sqrt{2}}, b = \sqrt{2}$ 中哪一个会是符合存在性的要求，但确定它们两个其中有一个会符合，所以是非构造性方法。事实上只要知道 $\sqrt{2}^{\sqrt{2}}$ 是无理数，令 $a = \sqrt{2}^{\sqrt{2}}, b = \sqrt{2}$ 这便是构造性方法。不过 $\sqrt{2}^{\sqrt{2}}$ 是无理数的证明非常困难（远远超过本讲义的范围）。所以我们避开它的证明，而仍能证明存在性，这就是非构造性方法的精神。

\begin{question}
试找到其他的例子或更一般的方法，利用构造性方法证明存在无理数 $a, b$使得 $a^{b}$为有理数。
\end{question}

大家可以发现利用构造性方法得到存在性，证明的重点并不是在于如何找到这些存在的东西，而是要确实验证并说明为何它们符合存在的条件。例如在证明一个方程式的解存在时，同学们常常利用等式推导出解的可能值，而没有代回验证是否为解就误以为证得存在性。这是同学们常弄错的，所以我们特别说明一下。在求解的推导过程，往往是假设解存在，再利用一些等式推导出解的 “可能值”。也就是说，除非你的推导过程是双向皆成立的，否则你推出的结果是，“若解存在的话，则它们可能的值”，未必这些真的是解（这再次说明若 $P$ 则 $Q$ 方向性的重要）。这些推出的值只是我们缩小要验证的范围而已，此时唯有将这些值代回验证，才能确保解的存在。我们看下面的例子。

\begin{example}
是否存在实数 $x$ 满足 $\sqrt{3 - 2x} = x - 2$？
\end{example}

若直接假设有解，则两边平方得 $3 - 2x = x^{2} - 4x + 4$，即 $x^{2} - 2x + 1 = 0$。得 $x = 1$。这表示若有解，其解必为 $x = 1$。也就是除了 $x = 1$ 可能会是此方程式之解外，其他实数都不可能是解。但将 $x = 1$ 代回原式，得 $1 = - 1$ 不合。故知不存在实数 $x$ 满足 $\sqrt{3 - 2x} = x - 2$。

其实使用非构造性方法证明存在性，一般或多或少会用到反证法。例如前面提过利用连续函数的中间值定理证明存在实数 $x$ 满足 $6x^{2} - x - 1 = 0$。虽然看似没用到反证法，不过中间值定理本身的证明一般来说都会用到反证法。还有一个常用来证明存在性的方法，就是鸽笼原理（pigeonhole principle）。\index{鸽笼原理（pigeonhole principle）}它原本称作“狄里赫利抽屉原理”（Dirichlet's drawer principle），不过现在许多人习惯用鸽笼原理称之。\index{狄里赫利抽屉原理（Dirichlet's drawer principle）|see{鸽笼原理（pigeonhole principle）}}

\begin{theorem}[鸽笼原理]
令 $n$ 为正整数。假设有 $n$ 个鸽笼以及多于 $n$ 只的鸽子。若要所有的鸽子住进鸽笼里，则一定会有一个鸽笼会有两只以上的鸽子。
\end{theorem}

\begin{proof}
很明显，这个存在性无法用构造性方法。利用反证法，原本命题 是说 “存在一个鸽笼会有两只以上的鸽子”，它的否定便是 “所有的鸽笼都只有一只或没有鸽子”。但如此一来表示所有 $n$ 个鸽笼里的鸽子最多只有 $n$ 只，与原本假设有多于 $n$ 只的鸽子矛盾。得证必有一个鸽笼会有两只以上的鸽子。
\end{proof}

将来我们会碰到利用鸽笼原理证明存在性的问题。大致上，只要弄清楚什么拿来当鸽子，什么拿来当鸽笼就好。例如证明任取 6 个整数，其中一定有两个整数其除以 5 的余数相同。这里我们可以将 6 个整数当成 6 只鸽子，将鸽笼从 0 到 4 编号。若除以 5 的余数为 0 就放到 0 号鸽笼，余数为 1 就放 1 号鸽笼，依此类推。因鸽子个数 6 多于鸽笼的个数 5，利用鸽笼原理，我们知道必有一个鸽笼里有两只以上的鸽子。也就是说会有两个整数除以 5 的余数相同。

我们可以将这个结果稍微推广一下。例如任取 16 个整数，可以证明其中可找到 4 个整数除以 5 有同样余数。这是因为若这些整数除以 5 余数为 0, 1, 2, 3, 4 的都不超过 3 个，那么这些整数总共最多不会超过 15 个，就和原本假设有 16 个整数相矛盾了。我们有定理2.2.3 的以下推广，证明就不再赘述了。

\begin{theorem}
令 $k, n$ 为正整数。假设有 $n$ 个鸽笼以及多于 $kn$ 只的鸽子。若要所有的鸽子住进鸽笼里，则一定会有一个鸽笼会有 $k$ 只以上的鸽子。
\end{theorem}

\begin{question}
试证明定理 2.2.4。
\end{question}

最后提醒一下，鸽笼原理不能处理鸽子数少于或等于鸽笼数的情况（除非另有特殊条件）。它并没有说在鸽子数少于或等于鸽笼数的情况下不会有鸽笼有两只以上的鸽子。另外它也没有说在鸽子数多于鸽笼数的情况下不会有空笼子。这些都很容易找到反例，请不要自己过度解读这个原理。

\subsection{唯一性}

基本上唯一性的证明是在假设存在的前提之下去证明唯一。所以唯一性的证明一般和存在性的证明是无关的。当然，如果已知不存在了，就不必去证明唯一性了。例如在 例 2.2.2 中当我们假设 $x$ 为解而推得 $x = 1$ 时，便证明了此方程式若有解则解唯一。不过后来知道方程式无解，所以这唯一性就不重要了。

大致上唯一性的证明也分成直接证明与反证法两种。直接证明就如前述，直接说明若东西存在应该是什么。而反证法一般用的方法是假设有两个不同的东西满足条件，进而推得矛盾。我们简单地用 $\mathbb{R}^{2}$ 上向量的性质来说明。

\begin{example}
证明 $\mathbb{R}^{2}$ 中若存在一个向量 $\overrightarrow{O}$ 满足对任意 $\mathbb{R}^{2}$ 上的向量 $\overrightarrow{V}$ 皆符合 $\overrightarrow{V} + \overrightarrow{O} = \overrightarrow{V}$，则 $\overrightarrow{O}$ 是唯一的。
\end{example}

(1) 直接证明：假设 $\overrightarrow{O} = (x, y)$，对任意 $\overrightarrow{V} = (a, b) \in \mathbb{R}^{2}$。由于 $\overrightarrow{O}$ 须符合 $\overrightarrow{V} + \overrightarrow{O} = \overrightarrow{V}$，得 $(a, b) + (x, y) = (a + x, b + y) = (a, b)$。利用向量相等的定义得 $a + x = a, b + y = b$，即 $x = 0, y = 0$。得证 $\overrightarrow{O}$ 若存在，则必须等于 $(0, 0)$。

(2) 反证法：假设 $\overrightarrow{O}, \overrightarrow{Q} \in \mathbb{R}^{2}$ 且 $\overrightarrow{O} \neq \overrightarrow{Q}$ 皆满足对任意 $\overrightarrow{V} \in \mathbb{R}^{2}$，
\begin{equation}
    \overrightarrow{V} + \overrightarrow{O} = \overrightarrow{V}
\end{equation}
以及
\begin{equation}
    \overrightarrow{V} + \overrightarrow{Q} = \overrightarrow{V}
\end{equation}
考虑 $\overrightarrow{Q} = \overrightarrow{V}$ 的情形代入式子 (2.1) 得 $\overrightarrow{Q} + \overrightarrow{O} = \overrightarrow{Q}$。同理将 $\overrightarrow{O} = \overrightarrow{V}$ 代入式子 (2.2) 得 $\overrightarrow{O} + \overrightarrow{Q} = \overrightarrow{O}$。由于 $\overrightarrow{Q} + \overrightarrow{O} = \overrightarrow{O} + \overrightarrow{Q}$，得 $\overrightarrow{Q} = \overrightarrow{O}$。此与当初假设 $\overrightarrow{O} \neq \overrightarrow{Q}$ 相矛盾，得证唯一性。

\begin{question}
给定 $\overrightarrow{V} \in \mathbb{R}^{2}$，试利用直接证法以及反证法证明：$\mathbb{R}^{2}$ 中若存在一个向量 $\overrightarrow{W}$ 满足 $\overrightarrow{V} + \overrightarrow{W} = \overrightarrow{O}$，则 $\overrightarrow{W}$ 是唯一的。
\end{question}

注意在 例 2.2.5 中的直接证法中，我们求出 $\overrightarrow{O}$ 若存在，则 $\overrightarrow{O} = (0, 0)$。如再代回验证，确认 $\overrightarrow{O} = (0, 0)$ 确实符合，我们便也证得存在性了。这是直接证法的好处。而反证法就没办法推得存在性了。所以一般不是用直接证明时，存在性及唯一性的证明是要分开来处理的。然而将来我们会碰到较抽象的数学问题时，大多直接证明是行不通的。此时只好仰赖反证法了。

\begin{example}
证明若存在一个实数 $r$ 满足 $r^3 = 3$，则此实数 $r$ 是唯一的。
\end{example}

\begin{proof}
回顾一下在 例 2.1.5 中我们证明了，若 $x, y$ 为实数且 $x \neq y$，则 $x^3 \neq y^3$。假设 $r \in \mathbb{R}$ 满足 $r^3 = 3$ 且 $s \neq r$ 是另一个实数满足 $s^3 = 3$，则利用 例 2.1.5 的结果得 $3 = s^3 \neq r^3 = 3$。由此矛盾知不可能有另一个实数 $s$ 会满足 $s^3 = 3$。
\end{proof}

在 例 2.2.6 中我们是先证明 $x \neq y$ 时，$x, y$ 不可能都符合某性质，再利用反证法推得唯一性。这是我们一般证明唯一性常用的方法。再强调一次，例 2.2.6 的证明中我们无法得知是否存在实数 $r$ 满足 $r^3 = 3$，我们只知道若存在的话必唯一。至于此存在性的证明，是需要另外利用实数的完备性（或利用多项式函数 $f(x) = x^3$ 的连续性）来证明的。这样有了存在性和唯一性我们才进一步用符号 $3^{1/3}$ 来表示这个唯一满足 $r^3 = 3$ 的实数 $r$。

\section{数学归纳法}

另一个常见的证明方法就是所谓的 “数学归纳法（mathematical induction）”\index{数学归纳法（mathematical induction）}。其实数学归纳法牵涉到建立整数系的公设（axiom），不过这里我们不去谈这些公设逻辑的问题。而着重于理解并正确使用数学归纳法。我们将介绍三种数学归纳法，虽然它们看起来不大一样，不过背后的原理是相同的，事实上它们是等价的。

数学的理论证明其实根源是一些大家都能接受但无法证明的公设。介绍数学归纳法之前，我们先了解所谓的\textbf{良序原理（well-ordering principle）}。\index{良序原理（well-ordering principle）}它是可以用其他的公设来证明的，不过由于我们目前不想牵涉到这方面的课题，所以我们直接把良序原理当成是一个公设。也就是它和我们的直观吻合，所以我们相信它而不去证明。

所谓良序，字面上的解释是 “好的排序” 的意思，这个排序原理简单来说是：“将一些正整数收集起来所成的非空集合中，一定有一个最小的元素”。相信大家应该不会觉得这个原则不对吧！直观上，自然数是有最小元素 1 的，所以有下界。因此大家应不会怀疑这集合里会有最小的元素。问题是这集合可能有无穷多的元素，我们可能无法真正找出这最小元素来，不过我们相信它一定存在。

接下来我们来看，最基本的第一种数学归纳法。

\begin{theorem}[数学归纳法]
假设以下两个命题是对的
\begin{enumerate}
    \item[(I1)] $P(1)$ 成立
    \item[(I2)] 若 $k \in \mathbb{N}$ 且 $P(k)$ 成立，则 $P(k+1)$ 成立
\end{enumerate}
那么对任意正整数 $n$，$P(n)$ 皆成立。
\end{theorem}

\begin{proof}
由于我们不可能直接证明所有的正整数 $n$ 都会使得 $P(n)$ 成立，所以我们用反证法。也就是假设 (1), (2) 是对的以及 “对任意正整数 $n$，$P(n)$ 皆成立” 是错的来得到矛盾。“对任意正整数 $n$，$P(n)$ 皆成立” 是错的，亦即 “存在正整数 $n$，使得 $P(n)$ 不成立”。因此我们可以将使得 $P(n)$ 不成立的这些正整数 $n$ 收集起来。因为它不是空集合，故由良序原理知，必存在最小的正整数 $m$ 使得 $P(m)$ 不成立。由 (1) 我们知道 $P(1)$ 成立，故得 $m \neq 1$。也就是说 $m$ 为大于 1 的正整数。现由于 $m - 1$ 为正整数且 $m - 1 < m$，故由 $m$ 为使得 $P(m)$ 不成立的最小正整数的假设知 $P(m - 1)$ 成立。然而由 (2) 知，当 $P(m - 1)$ 成立时，$P((m - 1) + 1) = P(m)$ 必成立。此与 $P(m)$ 不成立的假设相矛盾。故知不可能存在正整数 $n$，使得 $P(n)$ 不成立，也就是说对任意正整数 $n$, $P(n)$ 皆成立。
\end{proof}

数学归纳法是很好理解的，它是说由 (I1) 知 $P(1)$ 是对的，将 (I2) 代 $k = 1$ 的情况，故由 $P(1)$ 对可推得 $P(2)$ 是对的。接着代 $k = 2$ 的情况，由 $P(2)$ 是对的推得 $P(3)$ 是对的。这样一直下去。所以 $P(1)$ 的起头非常重要。另外要强调的是 (I2) 指的是假设 $P(k)$ 对推得 $P(k + 1)$ 是对的，所以它并不是要证明 $P(k)$ 是对的，你也不必担心 $P(k)$ 到底对不对，只要想法子利用 $P(k)$ 是对的假设证出 $P(k + 1)$ 是对的。若你没办法单纯由 $P(k)$ 是对的推得 $P(k + 1)$ 是对的，那基本上就无法用数学归纳法证明了。例如考虑多项式 $f(x) = x^{2} + x + 41$。当我们代 $x = 1$ 时 $f(1) = 43$ 为质数。代 $x = 2$，得 $f(2) = 47$ 仍为质数。$f(3) = 53$，$f(4) = 61$，$f(5) = 71$ 也都是质数。或许你会有一种冲动认为 $x$ 代任何的整数 $n$ 都会使 $f(n)$ 为质数。事实上一直代到 $x = 39$ 它都会是质数。但是光由代的动作，而不去考虑如何由 $f(k)$ 是质数得到 $f(k + 1)$ 是质数，是没有办法用数学归纳法的。事实上我们可以看出，当 $x = 40$ 时，$f(40) = 40^{2} + 40 + 41 = 40(40 + 1) + 41$ 是可以被 41 整除的，所以不是质数。

我们看下一个可以用数学归纳法证明的例子。

\begin{example}
设 $a, b$ 为相异的整数，利用数学归纳法证明：对任意的正整数 $n$ 皆有 $a^{n} - b^{n}$ 为 $a - b$ 的倍数。
\end{example}

\begin{proof}
我们可以将 $a^{n} - b^{n}$ 分解成 $(a - b)(a^{n - 1} + a^{n - 2}b + \dots + b^{n - 1})$，得证 $a^{n} - b^{n}$ 为 $a - b$ 的倍数。不过这里想介绍如何用数学归纳法证明。首先我们第一步代入 $n = 1$ 得 $a - b$，当然是 $a - b$ 的倍数，故成立。第二步是直接假设 $a^{k} - b^{k}$ 为 $a - b$ 的倍数，要推得 $a^{k + 1} - b^{k + 1}$ 为 $a - b$ 的倍数。所以我们要想办法看看 $a^{k + 1} - b^{k + 1}$ 和 $a^{k} - b^{k}$ 的关系。为了要让 $a^{k + 1} - b^{k + 1}$ 和 $a^{k} - b^{k}$ 扯上关系，我们自然将之写成
$$a^{k + 1} - b^{k + 1} = a a^{k} - b b^{k} = a a^{k} - a b^{k} + a b^{k} - b b^{k} = a(a^{k} - b^{k}) + (a - b)b^{k}.$$
此时由假设 $a^{k} - b^{k}$ 为 $a - b$ 的倍数，我们可将 $a^{k} - b^{k}$ 写成 $(a - b)m$，其中 $m$ 为整数。所以 $a^{k + 1} - b^{k + 1} = a(a - b)m + (a - b)b^{k} = (a - b)(am + b^{k})$。得证 $a^{k + 1} - b^{k + 1}$ 为 $a - b$ 的倍数。故由数学归纳法知，对任意的正整数 $n$，$a^{n} - b^{n}$ 为 $a - b$ 的倍数。
\end{proof}

一般来说，数学归纳法的第一步只是代值验证，应该没问题。而第二步就是一个若 $P$ 则 $Q$ 的证明，可以利用前面 2.1 节介绍的方法证明。在证明过程中若无法马上看出 $P(k)$ 和 $P(k + 1)$ 的关系，可以尝试先找出 $P(1), P(2)$ 的关系，$P(2), P(3)$ 的关系等，再推敲出 $P(k)$ 和 $P(k + 1)$ 的关系。另外要谨记，要单纯从 $P(k)$ 成立推出 $P(k + 1)$，在这样的过程中不能加入其他的假设。我们看一个错误的例子。

\begin{example}
以下的数学归纳法是要证明任意取 $n$ 个数都会相等。这个结论当然是错的，我们必须找出推论错误之处。

第一步：任取一个数，因为只有一个当然成立。

第二步：假设任取 $k$ 个数都会相等，要证明任取 $k + 1$ 个数也都会相等。现任取 $k + 1$ 个数，将这些数留下一个设其值为 $a$，而其他 $k$ 个数放入袋中。依假设袋中 $k$ 个数都会相等设为 $b$。现将袋中取出一个数，再将当初留下的那个数 $a$ 放入袋中。依假设此时袋中这 $k$ 个数都应相等，故得 $a = b$。因此依数学归纳法，得证任意取 $n$ 个数都会相等。
\end{example}

这个证明出错的当然是在第二步。它假设在袋中取出一个数后仍有 $k - 1$ 个数在袋中。然而当 $k = 1$ 时，此时袋中没有东西，最后放入袋中的数 $a$ 无其他的数与之比较，故无法得知 $a = b$。所以虽然上面的论述当 $k \ge 2$ 时，确实由假设 $P(k)$ 证得 $P(k + 1)$ 对，但在 $k = 1$ 时就无法由 $P(k)$ 对推得 $P(k + 1)$ 对。这不符合数学归纳法要求对所有的正整数 $k$，都要满足若 $P(k)$ 对，则 $P(k + 1)$ 对，所以论证是错误的。由这个例子我们建议，若要用数学归纳法，通常在验证 $P(1)$ 对之后，不要马上去证明 $P(k) \Rightarrow P(k + 1)$，而是想看看如何能单纯由 $P(1)$ 是对的推得 $P(2)$ 是对的。这样不只能让我们比较有想法知道如何去证明 $P(k) \Rightarrow P(k + 1)$，也可避免如上述的错误。

有时有的性质并不会从 $n = 1$ 开始就对，例如当 $n \ge 5$ 时，证明 $2^n > n^2$。这里数学归纳法若只能从 1 开始，就无法处理了。事实上，依照定理 2.3.1 的推论，我们有以下更一般化的数学归纳法。

\begin{corollary}[扩展数学归纳法，extended mathematical induction]\index{数学归纳法（mathematical induction）!扩展$\sim$（extended）}
设 $m$ 为整数。假设以下两个命题 是对的
\begin{enumerate}
    \item[(EI1)] $P(m)$ 成立
    \item[(EI2)] 若 $k \ge m$ 为整数且 $P(k)$ 成立，则 $P(k+1)$ 成立
\end{enumerate}
那么对任意大于等于 $m$ 的整数 $n$ 皆会使得 $P(n)$ 成立。
\end{corollary}

\begin{proof}
事实上，可以仿照定理2.3.1 用良序原理 来证明，不过这里我们想用 定理 2.3.1 来证明。首先令 $Q(n) = P(m + n - 1)$。因已知 $P(m)$ 成立，故知 $Q(1) = P(m)$ 成立。亦即 $Q$ 满足定理 2.3.1 的条件 (I1)。接着我们检查 $Q$ 是否符合定理 2.3.1 的条件 (I2)，也就是假设 $k \ge 1$ 为整数且 $Q(k)$ 成立，是否可推得 $Q(k+1)$ 成立。现假设 $k \in \mathbb{N}$ 且 $Q(k)$ 成立，此时 $m + k - 1 \ge m$ 为整数，且 $P(m + k - 1) = Q(k)$ 成立，故由 (EI2) 的假设得 $P(m + k) = Q(k + 1)$ 成立。我们证得，若 $k \in \mathbb{N}$ 且 $Q(k)$ 成立，则 $Q(k+1)$ 成立。故由 定理 2.3.1 知对所有的 $n^{\prime} \in \mathbb{N}$，$Q(n^{\prime}) = P(m + n^{\prime} - 1)$ 成立。因此当 $n$ 为大于等于 $m$ 的整数时，令 $n = m + n^{\prime} - 1$，此时 $n^{\prime} \in \mathbb{N}$，故得 $P(n) = P(m + n^{\prime} - 1) = Q(n^{\prime})$ 成立。
\end{proof}

这个 “扩展数学归纳法” 我们用 推论 称之，是因为它可由 定理 2.3.1 直接推得。事实上在 $m = 1$ 的情况 推论 2.3.4 就是 定理 2.3.1，所以我们知道 定理 2.3.1 和 推论 2.3.4 是等价的。

\begin{example}
证明对任意正整数 $n$，当 $n \ge 5$ 时，$2^n > n^2$。
\end{example}

\begin{proof}
我们用 扩展数学归纳法，$m = 5$ 且 $P(n)$ 为 $2^n > n^2$ 的情况证明。当 $n = 5$ 时，$2^5 = 32 > 25 = 5^2$，故 $P(5)$ 成立。假设 $k \ge 5$ 为整数且 $2^k > k^2$。因 $2^{k+1} = 2 \times 2^k$，故由 $2^k > k^2$ 之假设得 $2^{k+1} > 2k^2$。而 $(k+1)^2 = k^2 + 2k + 1$，若能证得 $2k^2 > k^2 + 2k + 1$，则得证 $2^{k+1} > (k+1)^2$，即 $P(k+1)$ 成立。然而 $2k^2 > k^2 + 2k + 1$ 等同于 $k^2 - 2k > 1$，又 $k^2 - 2k = k(k - 2)$，故由 $k \ge 5$ 得知 $k^2 - 2k > 1$。我们证得了若 $k \ge 5$ 为整数且 $P(k)$ 成立，则 $P(k+1)$ 成立，故由 扩展数学归纳法 (推论 2.3.4) 得证对任意大于等于 5 的整数 $n$ 皆会使得 $P(n)$，即 $2^n > n^2$ 成立。
\end{proof}

\begin{question}
在 例 2.3.5 的证明中要用到 $k(k - 2) > 1$。不过此式在 $k = 3$ 就会成立，为何 $2^n > n^2$ 需要到 $n \ge 5$ 时才会成立呢？
\end{question}

数学的理论推导，常常是由许多论证利用逻辑堆砌起来。因此在每一步骤，都应清楚这个步骤是在论证什么。常见的错误就是不明究理，盲目地推导，最后连自己的论证对错都不知道，甚至证出什么都不知道。若你能回答问题 2.8 这类的问题，那非常好，表示你能注意到每一步骤在证什么。在 例 2.3.5 我们是证明了当 $k \ge 3$ 时可由 $P(k)$ 成立推得 $P(k+1)$ 成立。它只告诉我们 $P(3)$ 成立的话 $P(4)$ 就会成立，并不表示 $P(3)$ 成立。事实上 $P(3), P(4)$ 都不成立。然而我们知道 $P(5)$ 成立，而又 $5 > 3$，所以可推得 $P(6)$ 成立，以至于对大于等于 5 的整数皆成立。同样的道理，假如 $P(5)$ 成立，而在推导 $P(k) \Rightarrow P(k+1)$ 的过程中发现只有在 $k \ge 10$ 的情况成立。此时无法推得 $P(6)$ 成立，所以也就无法马上下结论说 $P(n)$ 会对 $n \ge 5$ 都会成立了。但因为这里仅剩几个情况（即 $6 \le k \le 10$）需检验，若验证它们都成立，当然就能下结论说 $P(n)$ 会对 $n \ge 5$ 都会成立。

\begin{question}
假设当 $k \ge 10$ 皆会满足 $P(k) \Rightarrow P(k+1)$。试写下在以下情况当 $n$ 大于等于多少时，$P(n)$ 会成立（有可能无法确定）。
\begin{enumerate}
    \item $P(9)$ 成立。
    \item $P(11)$ 成立。
    \item $P(8), P(9), P(10)$ 成立。
    \item 当 $n$ 是 3 的倍数时 $P(n)$ 皆成立。
\end{enumerate}
\end{question}

在数学归纳法的证明中有时 $P(k)$ 对的条件不足以直接证明 $P(k+1)$ 对。例如一些递推数列的问题，有时需要之前更多项才能决定下一项的性质。事实上当我们由 $P(1)$ 对证得 $P(2)$ 对，再由 $P(2)$ 对要证 $P(3)$ 时，其实 $P(1)$ 已经知道是对的，所以我们不只有 $P(2)$ 对的前提，我们还有 $P(1)$ 对。同理在证得 $P(3)$ 对要证明 $P(4)$ 时，其实 $P(1), P(2)$ 为对的条件也可以用上，所以我们有以下条件更强的数学归纳法。

\begin{corollary}[强数学归纳法，strong mathematical induction]\index{数学归纳法（mathematical induction）!强$\sim$（strong）}
设 $m$ 为整数。假设以下两个命题 是对的
\begin{enumerate}
    \item[(SI1)] $P(m)$ 成立
    \item[(SI2)] 若 $k \ge m$ 为整数且 $P(m), P(m+1), \ldots , P(k-1), P(k)$ 皆成立，则 $P(k+1)$ 成立
\end{enumerate}
那么对任意大于等于 $m$ 的整数 $n$ 皆会使得 $P(n)$ 成立。
\end{corollary}

\begin{proof}
对于大于等于 $m$ 的整数 $n$，令 $Q(n) = P(m) \land P(m+1) \land \cdots \land P(n-1) \land P(n)$。因假设 $P(m)$ 成立，故知 $Q(m) = P(m)$ 成立，亦即 $Q$ 满足 推论 2.3.4 的条件 (EI1)。接着我们检查 $Q$ 是否符合 推论 2.3.4 的条件 (EI2)，也就是假设 $k \ge m$ 为整数且 $Q(k)$ 成立，是否可推得 $Q(k+1)$ 成立。然而 $Q(k)$ 成立表示 $P(m), \ldots , P(k-1), P(k)$ 皆成立，故由 (SI2) 的假设得 $P(k+1)$ 成立。然而已知 $P(m), P(m+1), \ldots , P(k-1), P(k)$ 皆成立，故知 $Q(k+1) = P(m) \land P(m+1) \land \cdots \land P(k) \land P(k+1)$ 成立。我们证得，若 $k \ge m$ 为整数且 $Q(k)$ 成立，则 $Q(k+1)$ 成立。故由 推论 2.3.4 知对任意大于等于 $m$ 的整数 $n$ 皆会使得 $Q(n)$ 成立。然而 $Q(n)$ 成立表示 $P(m), P(m+1), \ldots , P(n-1), P(n)$ 皆成立，自然 $P(n)$ 成立，故得证当 $n$ 为大于等于 $m$ 的整数时，$P(n)$ 皆成立。
\end{proof}

\begin{remark}
强数学归纳法 的 (SI2) 一般我们可以改写为：若 $k \ge m$ 为整数且对于满足 $m \le i \le k$ 的整数 $i$, $P(i)$ 皆成立，则 $P(k+1)$ 成立。
\end{remark}

我们称 推论 2.3.6 为 强数学归纳法，意即它比 定理 2.3.1 强。在数学上，我们称一个定理比另一个定理强，大致上的意思是它可以在更广泛的情况使用。例如 推论 2.3.4 就比 定理 2.3.1 强，因为它可以应用在任意整数 $m$ 起头的情况，而不只是 1。通常比较强的定理很容易就可推导出比较弱的。例如 推论 2.3.4 只要考虑 $m = 1$ 的情况就可得 定理 2.3.1。不过感觉上比较弱的定理，未必就真的比较弱。例如我们是用 定理 2.3.1 证得 推论 2.3.4，所以逻辑上来说他们是等价的，没有谁强谁弱之分。不过在使用上当然 推论 2.3.4 比较方便。

同样地，推论 2.3.6 直观上也比 推论 2.3.4 强（因为 (SI2) 的条件比较多、比较好使用）。我们可以很容易用 推论 2.3.6 来证明 推论 2.3.4。事实上当 $P$ 符合 推论 2.3.4 的 (EI1), (EI2)，我们希望证明 $P$ 也会符合 推论 2.3.6 的 (SI1), (SI2)，因此由 推论 2.3.6 的结论得到 $P(n)$ 对所有 $n \ge m$ 都成立。其实 (EI1) 和 (SI1) 是一样的，所以只要探讨 $P$ 若符合 (EI2) 是否会符合 (SI2)，也就是假设 $P(m), \ldots , P(k - 1), P(k)$ 成立是否可得 $P(k + 1)$ 成立。然而由这个假设当然可得 $P(k)$ 成立，而又已知 $P$ 符合 (EI2)，也就是 $P(k)$ 成立的话 $P(k + 1)$ 一定成立，也因此推得 $P$ 确实符合 (SI2)。另一方面，推论 2.3.6 的证明是利用 推论 2.3.4 证得的，也因此知它们是等价的。最后由等价的传递性，我们知道里介绍的三个数学归纳法都是等价的。这是从逻辑的观点来看（它们没有强弱之分），实际在证明时当然是挑最适合的来证明。下一个例子我们可以看出，应该用 强数学归纳法 来证明较合适。

\begin{example}
证明所有大于 1 的整数都可以写成有限多个质数的乘积。
\end{example}

或许很多同学会用以下的方法证明。设 $n$ 为大于 1 的整数，若 $n$ 为质数，则 $n$ 符合可写成有限多个质数乘积。若 $n$ 不是质数，依定义 $n$ 可以写成两个比 $n$ 小但大于 1 的整数相乘，这样一直下去可证得可以写成有限多个质数的乘积。相信大家都能接受这样的说法来解释这个性质是对的。但它并不是好的证明，比方说如何 “一直下去” 且为何这个程序经有限多次后会停止（这样才能说是有限多个质数的乘积）也要说明清楚。有了数学归纳法，就是能帮我们解决这些不容易说清楚的地方。在这个例子，若我们仅假设 $k$ 可以写成有限多个质数乘积，是无法证得 $k + 1$ 可以写成有限多个质数乘积，所以我们必须用 强数学归纳法 来证明。

\begin{proof}
当 $n = 2$ 时，因 2 为质数，故成立。假设当 $k \ge 2$ 时所有满足 $2 \le i \le k$ 的整数 $i$ 都可以写成有限多个质数的乘积。现考虑 $k + 1$ 的情形。因为 $k + 1$ 是质数时自然成立，所以我们仅需考虑 $k + 1$ 不为质数的情形。此时 $k + 1 = ab$，其中 $1 < a, b \le k$。故由前归纳之假设知 $a, b$ 皆为有限多个质数的乘积，因此 $k + 1 = ab$ 自然可以写成有限多个质数的乘积。故由 强数学归纳法 知所有大于 1 的整数都可以写成有限多个质数的乘积。
\end{proof}

在利用数学归纳法证明的过程中，最重要且最难的部分就是第二步骤——由假设 $P(k)$ 成立（或 $P(m), P(m + 1), \ldots , P(k)$ 成立）证得 $P(k + 1)$ 成立。这里常常在推导的过程中发现 $k$ 要在某些范围内才会对。例如在 例 2.3.3 的错误示范中，其实 “任取 $k$ 个数相等推得任取 $k + 1$ 个数会相等” 的证明在 $k \ge 2$ 时是对的。所以此时若能再补上 $k = 1$ 的情况也对，整个证明就完成了（当然在 例 2.3.3 这个例子这是不可能做到的）。前面提过，当我们在处理第二步骤时，若发现 $k$ 要有所限制才能对，此时我们可以多检查那些 $k$ 无法涵盖的情况，若这些情况也都对，就完成归纳法的证明了。总之，在数学归纳法的证明中，有时并不是仅检查初始的情况就好，我们再多看一些例子。

\begin{example}
考虑所谓的斐波那契数列（Fibonacci sequence） $\{F_0, F_1, F_2, \ldots \}$，即 $F_0 = 0, F_1 = 1$ 且对任意 $i \ge 2$, $F_i$ 满足 $F_i = F_{i - 1} + F_{i - 2}$。证明$n \ge 4$时 $F_n < 2^{n - 2}$。
\end{example}

很显然地，因 $F_{k + 1}$ 的值由 $F_{k}$ 和 $F_{k - 1}$ 所决定，我们无法仅由 $F_{k}$ 来推得 $F_{k + 1}$。所以这里我们用 强数学归纳法 来处理。依定义 $F_2 = F_1 + F_0 = 1$, $F_3 = F_2 + F_1 = 1 + 1 = 2$，故当 $n = 4$ 时，$F_4 = F_3 + F_2 = 2 + 1 = 3$，所以 $F_4 = 3 < 2^{4 - 2} = 4$ 成立。现假设 $k \ge 4$ 且对所有 $4 \le i \le k$，皆有 $F_i < 2^{i - 2}$。此时 $F_{k + 1} = F_k + F_{k - 1}$。我们希望用到 $F_k < 2^{k - 2}$ 和 $F_{k - 1} < 2^{(k - 1) - 2} = 2^{k - 3}$ 的假设推得 $F_{k + 1} < 2^{(k + 1) - 2} = 2^{k - 1}$。不过当 $k = 4$ 时，$i = k - 1$ 并不符合 $4 \le i \le k$，所以此时无法使用 $F_{k - 1} < 2^{k - 3}$ 的假设（事实上此时 $F_{k - 1} = F_3 = 2 = 2^{4 - 3}$）。所以我们再补上 $k = 4$ 的情况，即直接验证 $F_{k + 1} = F_5 = F_4 + F_3 = 5 < 2^{5 - 2} = 8$，才可完成证明。

\begin{proof}
首先直接验证得 $F_4 = 3 < 2^{4 - 2}$, $F_5 = 5 < 2^{5 - 2}$。

现假设 $k \ge 5$ 且对任意 $i = 4, 5, \ldots , k$ 皆有 $F_i < 2^{i - 2}$。因为 $4 \le k - 1 \le k$ 且 $4 \le k \le k$，我们有 $F_k < 2^{k - 2}$, $F_{k - 1} < 2^{(k - 1) - 2} = 2^{k - 3}$，故得
\[
\begin{aligned}
    F_{k + 1} = F_{k} + F_{k - 1} &< 2^{k - 2} + 2^{k - 3} \\&= 2^{k - 3}(2 + 1) \\&< 4 \times 2^{k - 3} = 2^{k - 1} = 2^{(k + 1) - 2}.
\end{aligned}
\]
依数学归纳法得证$n \ge 4$时 $F_n < 2^{n - 2}$。
\end{proof}

在下一个例子中，我们想证明所有大于 20 的整数都可以写成 4 和 5 的正整数倍之和。也就是证明若 $n > 20$，则存在正整数 $l, m$ 使得 $n = 4l + 5m$。或许同学会想到 $1 = 5 - 4$，所以若 $k = 4l + 5m$，则 $k + 1 = 4l + 5m + (5 - 4) = 4(l - 1) + 5(m + 1)$。不错，用这个方法可以证明所有的整数皆可以写成 4 的倍数和 5 的倍数之和，但我们要求的是写成 4 和 5 的正整数倍之和。因此无法用这方法证明。不过观察一下发现若 $k = 4l + 5m$，则 $k + 4n = 4(l + n) + 5m$，所以我们可以先验证 21, 22, 23, 24 都可以写成 4 和 5 的正整数倍之和，就可利用 分类讨论将所有大于 20 的正整数分成 $21 + 4n$, $22 + 4n$, $23 + 4n$, $24 + 4n$ 来探讨，而得证。根据这个想法，我们用 强数学归纳法 来证明。

\begin{example}
证明若 $n$ 为整数且 $n > 20$，则存在 $l, m$ 为正整数满足 $n = 4l + 5m$。
\end{example}

\begin{proof}
由于 $21 = 4 \times 4 + 5 \times 1$, $22 = 4 \times 3 + 5 \times 2$, $23 = 4 \times 2 + 5 \times 3$ 和 $24 = 4 \times 1 + 5 \times 4$，故知当 $n = 21, 22, 23, 24$ 时成立。现假设 $k \ge 24$ 时对于所有满足 $21 \le i \le k$ 的整数 $i$，皆存在正整数 $l, m$ 使得 $i = 4l + 5m$。考虑 $k + 1$ 的情形，由于 $k + 1 = (k - 3) + 4$，且 $i = k - 3$ 满足 $21 \le i \le k$，故存在正整数 $l, m$ 使得 $k - 3 = 4l + 5m$，得 $k + 1 = 4l + 5m + 4 = 4(l + 1) + 5m$。由数学归纳法知，当 $n$ 为大于 20 的整数，存在 $l, m$ 为正整数满足 $n = 4l + 5m$。
\end{proof}

数学归纳法是一个很好使用的数学工具。它不只可以拿来处理整数的问题，其实很多可以用整数分类的问题都可以用数学归纳法处理。例如代数中许多有关多项式的问题，我们都可以用数学归纳法处理。

\begin{problemset}
    \item 使用指定的步骤直接证明：如果 $n$ 是整数，那么 $n^2 + 3n + 2$ 是偶数。
    \begin{probenumerate}
        \item 因式分解 $n^2 + 3n + 2$。
        \item 将 $n$ 分为两种情况。
    \end{probenumerate}

    \item 假设 $a, b$ 是整数。令 $P$ 为命题：如果 $a \times b$ 是偶数，那么 $a$ 或 $b$ 至少有一个是偶数。
    \begin{probenumerate}
        \item 写出 $P$ 的逆否命题。
        \item 利用直接证明法证明 $P$（提示：你可以假设 $a$ 是奇数）。
        \item 利用逆否命题法证明 $P$。
    \end{probenumerate}

    \item 令 $Q$ 为命题：如果 $a, b$ 是正整数，那么 $a^2 - b^2 \neq 1$。
    \begin{probenumerate}
        \item 否定 $Q$。
        \item 利用分情况证明法证明 $Q$（提示：(1) $a < b$; (2) $a = b$; (3) $a > b$，并且对于 (3) 设 $a = b + k$）。
        \item 利用反证法证明 $Q$。
    \end{probenumerate}

    \item 平面上有 5 个点，不全共线。证明存在一条直线仅通过其中两个点。

    \item 给定 $a, b \in \mathbb{N}$ 且 $b > 1$，证明如果存在 $h, r \in \mathbb{Z}$ 满足 $0 \leq r < b$ 使得 $a = bh + r$，那么 $h$ 是唯一的且 $r$ 是唯一的。

    \item 利用鸽笼原理证明：无论如何在 $8 \times 9$ 的长方形内放置 7 个点，总会有一对点的距离最大为 5。（提示：$3^2 + 4^2 = 5^2$）

    \item 利用数学归纳法证明鸽笼原理。

    \item 利用（强）数学归纳法证明以下内容：
    \begin{probenumerate}
        \item 如果 $a_1 = 1$ 且当 $n \geq 2$ 时 $a_n = 2a_{n-1} + 1$，那么对所有 $n \in \mathbb{N}$，$a_n = 2^n - 1$。
        \item 如果 $b_1 = 1$ 且当 $n \geq 2$ 时 $b_n = n + \sum_{i=1}^{n-1} b_i$，那么对所有 $n \in \mathbb{N}$，$b_n = 2^n - 1$。
    \end{probenumerate}
\end{problemset}
